\section{从归纳总结到物理世界 AI}

高继扬反复讲一个方法：归纳总结。小学、物理竞赛、博士发论文、工作面试和创业选择，都被他处理成“先设目标，再拆路径，再提高命中概率”。这种方法听起来朴素，却解释了他为什么会从清华电子系、深度学习、计算机视觉、自动驾驶一路走到具身智能。

\begin{figure}[H]
\centering
\includegraphics[width=0.96\textwidth]{trajectory-industrial-training.png}
\caption{高继扬的产业训练路径：归纳总结、AI 入门、Waymo、Momenta、星海图。}
\end{figure}

\begin{knowledgebox}{读图：三种能力如何叠加}
图中路线不是履历罗列。早期物理竞赛和博士阶段训练的是归纳、拆解和概率管理；商汤实习让他第一次感到神经网络能从数据中自动提炼规律；Waymo 训练系统工程和测量；Momenta 训练产品交付和客户价值；星海图则把这些能力带到具身智能。
\end{knowledgebox}

\subsection{曾国藩：从清流到事功}

前面讲高继扬的“归纳总结”，这一节要解释他的现实主义从哪里来。曾国藩这个例子并不是历史装饰，而是他理解创业组织的一个早期模板。

访谈中高继扬谈到申请学校受挫后读曾国藩。他关注的不是历史趣味，而是“一个儒家清流怎样变成能调动资源、组织队伍、打现实世界硬仗的人”。这段内容是理解他创业观的钥匙：做现实世界的事，不能只靠理念、论文或聪明，而要能拉动资源、组织人、承担结果。

\begin{importantbox}{事功视角}
所谓事功，不是短期功利，而是把目标落实成资源、组织、行动和结果。机器人创业尤其需要这种视角，因为它的每一步都要在现实世界里兑现。
\end{importantbox}

\subsection{深度学习的吸引力}

在商汤实习时，他第一次训练神经网络，感受到“机器从数据里自己提炼规律”的魔力。传统编程是人总结规则，再把规则写成代码；机器学习则把规则压进参数，让模型从数据分布中学习。这一刻奠定了他后面对 AI 的理解：AI 的价值不只是自动化，而是改变人类发现规律和构造生产力的方式。

\begin{knowledgebox}{术语消化：从规则到学习}
\begin{tabularx}{\textwidth}{p{0.22\textwidth}p{0.32\textwidth}X}
\toprule
术语 & 解决的问题 & 本期中的意义 \\
\midrule
Rule-based & 人手写规则、模块和 if-else & 传统自动驾驶和机器人系统的重要底座。 \\
Data-driven & 从数据中学习规律 & 高继扬选择 AI 和自动驾驶的核心理由。 \\
Benchmark & 用统一评测衡量模型表现 & 发论文、模型迭代和数据 recipe 都需要它。 \\
Physical World AI & AI 直接作用于物理世界任务 & 自动驾驶是第一种形态，具身智能是下一种形态。 \\
\bottomrule
\end{tabularx}
\end{knowledgebox}

\subsection{本章小结}

高继扬的主线是：把个人勤奋转化成可复用方法，再把 AI 的数据驱动能力放到物理世界里。自动驾驶和机器人吸引他的地方，正是 AI 成为行业底层变量，而不是局部优化工具。

